\section{Hamilton paths with prescribed colours}\label{sec:hamilton-paths}

\topic{Hahn's question}

\begin{problem}[Hahn, 1980, as quoted by Maamoun--Meyniel~\cite{mm84}]\label{prob:hahn}
Let $n\geq2k$, $k\geq2$ and $(n,k)\neq(4,2)$. In every colouring of the edges of $K_n$ in which each colour appears at most $k$ times, is there a rainbow Hamilton path?
\end{problem}

The intuition is that few edges per colour should make a rainbow path easy to find. A colouring is \term{proper} if adjacent edges receive different colours; then each colour class is a matching, of size at most $n/2$. An algebraic colouring answers Hahn's question in the negative.

\begin{theorem}[Maamoun--Meyniel~\cite{mm84}]\label{thm:mm}
Let $n=2^p$ with $p\geq2$, identify $V(K_n)$ with $\mathbb F_2^p$, and give the edge $xy$ the colour $x+y$. This colouring is proper, every colour appears exactly $n/2$ times, and there is no rainbow Hamilton path.
\end{theorem}

\begin{proof}
A rainbow Hamilton path $x_1,\dots,x_n$ has $n-1$ edges, so it uses each non-zero colour exactly once. In $\mathbb F_2^p$,
\[
x_n-x_1=\sum_{i}(x_{i+1}-x_i)=\sum_{c\neq0}c=0,
\]
so $x_n=x_1$, a contradiction.
\end{proof}

\begin{figure}[ht]
\centering
\begin{tikzpicture}[scale=1.05, every node/.style={font=\footnotesize}]
    \node[overtex] (a) at (0,1.5) {}; \node[above left=1pt of a] {$00$};
    \node[overtex] (b) at (1.7,1.5) {}; \node[above right=1pt of b] {$01$};
    \node[overtex] (c) at (1.7,0) {}; \node[below right=1pt of c] {$11$};
    \node[overtex] (d) at (0,0) {}; \node[below left=1pt of d] {$10$};
    \draw[col1] (a)--(b) node[midway,above] {$01$};
    \draw[col1] (d)--(c) node[midway,below] {$01$};
    \draw[col2] (a)--(d) node[midway,left] {$10$};
    \draw[col2] (b)--(c) node[midway,right] {$10$};
    \draw[col3] (a)--(c) node[pos=.42,above right] {$11$};
    \draw[col3] (b)--(d) node[pos=.42,below right] {$11$};
  \end{tikzpicture}

\caption{The sum colouring of $K_{\mathbb F_2^2}$: each colour is a perfect matching, and none of the $24$ orderings of the vertices is rainbow.}
\label{fig:maamoun-meyniel}
\end{figure}

\topic{Buratti's problem}

For $x,y\in\Z_v$ the \term{length} of the edge $xy$ of $K_v$ is $\ell(x,y)=\min(|x-y|,\,v-|x-y|)$~\cite[\S1]{hr09}.

\begin{conjecture}[Buratti, 2007, communicated to Rosa~\cite{hr09}]\label{conj:buratti}
Let $p=2n+1$ be prime and let $L$ be any multiset of $2n$ elements of $\{1,\dots,n\}$. Then $K_p$ has a Hamilton path whose multiset of edge lengths is exactly $L$.
\end{conjecture}

\begin{figure}[ht]
\centering
\begin{tikzpicture}[every node/.style={font=\footnotesize}]
    \polygonPoints{q}{7}{1.3}{0}
    \foreach \i in {0,...,6} \foreach \j in {0,...,6} {
      \ifnum\i<\j \draw[gray!20] (q\i)--(q\j); \fi}
    \draw[col1, line width=1.3pt] (q0)--(q1);
    \draw[col1, line width=1.3pt] (q1)--(q2);
    \draw[col3, line width=1.3pt] (q2)--(q5);
    \draw[col2, line width=1.3pt] (q5)--(q3);
    \draw[col3, line width=1.3pt] (q3)--(q6);
    \draw[col2, line width=1.3pt] (q6)--(q4);
    \polygonLabels{q}{7}{1.3}{0}
  \end{tikzpicture}

\caption{$p=7$ and $L=\{1,1,2,2,3,3\}$: the Hamilton path $0,1,2,5,3,6,4$.}
\label{fig:buratti}
\end{figure}

With only one length the conjecture is easy, and with two lengths it was solved by Horak and Rosa~\cite{hr09}. Meszka verified it by computer for all $p\leq23$~\cite[\S1]{hr09}; for $p=7$ all $28$ lists are realised. Horak and Rosa extended the question to all $v$.

\begin{conjecture}[Horak--Rosa~\cite{hr09}; in the form of Pasotti--Pellegrini~\cite{pp14}]\label{conj:bhr}
A multiset $L$ of $v-1$ elements of $\{1,\dots,\lfloor v/2\rfloor\}$ is the list of edge lengths of a Hamilton path of $K_v$ if and only if, for every divisor $d$ of $v$, at most $v-d$ elements of $L$ are multiples of $d$.
\end{conjecture}

The condition is necessary. An edge whose length is a multiple of $d$ joins two vertices in the same coset of the subgroup $d\Z_v$, which has $d$ cosets. The edges of the path with such lengths form a forest whose components stay inside cosets, so it has at least $d$ components and hence at most $v-d$ edges. For example, for $v=4$ and $L=\{2,2,2\}$ there are three multiples of $2$ but $v-d=2$; indeed only the two edges $\{0,2\}$ and $\{1,3\}$ have length $2$. For $v=p$ prime the condition is empty, and Conjecture~\ref{conj:bhr} is Buratti's conjecture.
